% GENERATED FILE. Edit canonical lessons, then run npm run generate:books.
% canonical-source-hash: fnv1a64:8bc6a6606a916376
% canonical-lessons: ML-C150-urulakkilannu, ML-C150-takkali, ML-C150-valutana, ML-C150-kurumulaku, ML-C150-veluttulli, ML-R150-first-pass-can-want-and-why, ML-R150-second-pass-why-and-the-house

\chapter{Potato, Tomato, Aubergine, Black pepper, Garlic}
\label{ch:ml-potato-tomato-aubergine-black-pepper-garlic}
\hlchaptermodality{150}

\begin{chapteropening}
\textbf{By the end of this chapter:} \emph{I can say a potato, a tomato, an aubergine, black pepper and garlic.}

Complete the last lesson of chapter 150: I can say a potato, a tomato, an aubergine, black pepper and garlic.
\end{chapteropening}

\section[\texorpdfstring{\ml{ഉരുളക്കിഴങ്ങ്} (urulakkiḻaṅṅŭ)}{ഉരുളക്കിഴങ്ങ് (urulakkiḻaṅṅŭ)}]{\ml{ഉരുളക്കിഴങ്ങ്} (urulakkiḻaṅṅŭ) — a potato}

\label{lesson:ML-C150-urulakkilannu}



\begin{quote}
Before the new one: say the Malayalam for a banana, then the Malayalam for grapes.
\end{quote}

\subsection*{You'll want to know: \ml{ഉരുളക്കിഴങ്ങ്}}
\textbf{\ml{ഉരുളക്കിഴങ്ങ്}} — \emph{urulakkiḻaṅṅŭ} — \textquotedblleft{}a potato\textquotedblright{}.

\begin{grammarlens}[title={What you've built}]
One more word for this chapter.
\end{grammarlens}

\subsection*{Your turn}
\begin{itemize}
\raggedright
  \item \emph{Say it:} \emph{urulakkiḻaṅṅŭ}
  \item \emph{Say it:} \emph{urulakkiḻaṅṅŭ}, once more
  \item \emph{Say it:} say \emph{urulakkiḻaṅṅŭ}
  \item \emph{Recall:} say the Malayalam for a banana, then the Malayalam for grapes, then say \emph{urulakkiḻaṅṅŭ} again
\end{itemize}

\begin{tcolorbox}[breakable,colback=teal!4,colframe=teal!35!black,title={Before you move on}]
What does \ml{ഉരുളക്കിഴങ്ങ്} mean? (\textquotedblleft{}A potato\textquotedblright{}.) Say it once more.
\end{tcolorbox}
\section[\texorpdfstring{\ml{തക്കാളി} (takkāḷi)}{തക്കാളി (takkāḷi)}]{\ml{തക്കാളി} (takkāḷi) — a tomato}

\label{lesson:ML-C150-takkali}



\begin{quote}
Before the new one: say the Malayalam for grapes, then the Malayalam for a potato.
\end{quote}

\subsection*{You'll want to know: \ml{തക്കാളി}}
\textbf{\ml{തക്കാളി}} — \emph{takkāḷi} — \textquotedblleft{}a tomato\textquotedblright{}.

\begin{grammarlens}[title={What you've built}]
One more word for this chapter.
\end{grammarlens}

\subsection*{Your turn}
\begin{itemize}
\raggedright
  \item \emph{Say it:} \emph{takkāḷi}
  \item \emph{Say it:} \emph{takkāḷi}, once more
  \item \emph{Say it:} say \emph{takkāḷi}
  \item \emph{Recall:} say the Malayalam for grapes, then the Malayalam for a potato, then say \emph{takkāḷi} again
\end{itemize}

\begin{tcolorbox}[breakable,colback=teal!4,colframe=teal!35!black,title={Before you move on}]
What does \ml{തക്കാളി} mean? (\textquotedblleft{}A tomato\textquotedblright{}.) Say it once more.
\end{tcolorbox}
\section[\texorpdfstring{\ml{വഴുതന} (vaḻutana)}{വഴുതന (vaḻutana)}]{\ml{വഴുതന} (vaḻutana) — an aubergine}

\label{lesson:ML-C150-valutana}



\begin{quote}
Before the new one: say the Malayalam for a potato, then the Malayalam for a tomato.

\emph{Uccakaḻiññŭ} is two words run together. Which one means \textquotedblleft{}noon\textquotedblright{}? (\emph{Ucca}.)
\end{quote}

\subsection*{You'll want to know: \ml{വഴുതന}}
\textbf{\ml{വഴുതന}} — \emph{vaḻutana} — \textquotedblleft{}an aubergine\textquotedblright{}.

\begin{grammarlens}[title={What you've built}]
One more word for this chapter.
\end{grammarlens}

\subsection*{Your turn}
\begin{itemize}
\raggedright
  \item \emph{Say it:} \emph{vaḻutana}
  \item \emph{Say it:} \emph{vaḻutana}, once more
  \item \emph{Say it:} say \emph{vaḻutana}
  \item \emph{Recall:} say the Malayalam for a potato, then the Malayalam for a tomato, then say \emph{vaḻutana} again
\end{itemize}

\begin{tcolorbox}[breakable,colback=teal!4,colframe=teal!35!black,title={Before you move on}]
What does \ml{വഴുതന} mean? (\textquotedblleft{}An aubergine\textquotedblright{}.) Say it once more.
\end{tcolorbox}
\section[\texorpdfstring{\ml{കുരുമുളക്} (kurumuḷakŭ)}{കുരുമുളക് (kurumuḷakŭ)}]{\ml{കുരുമുളക്} (kurumuḷakŭ) — black pepper}

\label{lesson:ML-C150-kurumulaku}



\begin{quote}
Before the new one: say the Malayalam for a tomato, then the Malayalam for an aubergine.
\end{quote}

\subsection*{You'll want to know: \ml{കുരുമുളക്}}
\textbf{\ml{കുരുമുളക്}} — \emph{kurumuḷakŭ} — \textquotedblleft{}black pepper\textquotedblright{}.

\begin{grammarlens}[title={What you've built}]
One more word for this chapter.
\end{grammarlens}

\subsection*{Your turn}
\begin{itemize}
\raggedright
  \item \emph{Say it:} \emph{kurumuḷakŭ}
  \item \emph{Say it:} \emph{kurumuḷakŭ}, once more
  \item \emph{Say it:} say \emph{kurumuḷakŭ}
  \item \emph{Recall:} say the Malayalam for a tomato, then the Malayalam for an aubergine, then say \emph{kurumuḷakŭ} again
\end{itemize}

\begin{tcolorbox}[breakable,colback=teal!4,colframe=teal!35!black,title={Before you move on}]
What does \ml{കുരുമുളക്} mean? (\textquotedblleft{}Black pepper\textquotedblright{}.) Say it once more.
\end{tcolorbox}
\section[\texorpdfstring{\ml{വെളുത്തുള്ളി} (veḷuttuḷḷi)}{വെളുത്തുള്ളി (veḷuttuḷḷi)}]{\ml{വെളുത്തുള്ളി} (veḷuttuḷḷi) — garlic}

\label{lesson:ML-C150-veluttulli}



\begin{quote}
Before the new one: say the Malayalam for an aubergine, then the Malayalam for black pepper.

Telugu widens its word for noon to cover the afternoon. What does Malayalam do instead? (It builds a phrase, \emph{uccakaḻiññŭ}.)
\end{quote}

\subsection*{You'll want to know: \ml{വെളുത്തുള്ളി}}
\textbf{\ml{വെളുത്തുള്ളി}} — \emph{veḷuttuḷḷi} — \textquotedblleft{}garlic\textquotedblright{}.

\begin{grammarlens}[title={What you've built}]
One more word for this chapter.
\end{grammarlens}

\subsection*{Your turn}
\begin{itemize}
\raggedright
  \item \emph{Say it:} \emph{veḷuttuḷḷi}
  \item \emph{Say it:} \emph{veḷuttuḷḷi}, once more
  \item \emph{Say it:} say \emph{veḷuttuḷḷi}
  \item \emph{Recall:} say the Malayalam for an aubergine, then the Malayalam for black pepper, then say \emph{veḷuttuḷḷi} again
\end{itemize}

\begin{tcolorbox}[breakable,colback=teal!4,colframe=teal!35!black,title={Before you move on}]
What does \ml{വെളുത്തുള്ളി} mean? (\textquotedblleft{}Garlic\textquotedblright{}.) Say it once more.
\end{tcolorbox}
\section[(dialogue)]{Can, want and why — a first pass}

\label{lesson:ML-R150-first-pass-can-want-and-why}



\begin{quote}
Say the words for black pepper and garlic.
\end{quote}

\subsection*{Your turn}
Say these aloud:

\begin{itemize}
\raggedright
  \item ability, strength, a right, an opportunity, a tool
  \item wealth, an owner, experience, a habit, a qualification
  \item so, a purpose, a result, a cause, on account of
  \item a wish, a dream, hope, an aim, a plan
\end{itemize}

\begin{tcolorbox}[breakable,colback=teal!4,colframe=teal!35!black,title={Before you move on}]
Ability? (\textbf{\ml{കഴിവ്}}, \emph{kaḻivŭ}.) Strength? (\textbf{\ml{ശക്തി}}, \emph{śakti}.) A right? (\textbf{\ml{അവകാശം}}, \emph{avakāśaṁ}.) An opportunity? (\textbf{\ml{അവസരം}}, \emph{avasaraṁ}.) A tool? (\textbf{\ml{ഉപകരണം}}, \emph{upakaraṇaṁ}.) Wealth? (\textbf{\ml{സമ്പത്ത്}}, \emph{sampattŭ}.) An owner? (\textbf{\ml{ഉടമസ്ഥൻ}}, \emph{uṭamastan}.) Experience? (\textbf{\ml{പരിചയം}}, \emph{paricayaṁ}.) A habit? (\textbf{\ml{ശീലം}}, \emph{śīlaṁ}.) A qualification? (\textbf{\ml{യോഗ്യത}}, \emph{yōgyata}.) So? (\textbf{\ml{അതുകൊണ്ട്}}, \emph{atukoṇṭŭ}.) A purpose? (\textbf{\ml{ഉദ്ദേശ്യം}}, \emph{uddēśyaṁ}.) A result? (\textbf{\ml{ഫലം}}, \emph{phalaṁ}.) A cause? (\textbf{\ml{ഹേതു}}, \emph{hētu}.) On account of? (\textbf{\ml{നിമിത്തം}}, \emph{nimittaṁ}.) A wish? (\textbf{\ml{ആഗ്രഹം}}, \emph{āgrahaṁ}.) A dream? (\textbf{\ml{സ്വപ്നം}}, \emph{svapnaṁ}.) Hope? (\textbf{\ml{പ്രതീക്ഷ}}, \emph{pratīkṣa}.) An aim? (\textbf{\ml{ലക്ഷ്യം}}, \emph{lakṣyaṁ}.) A plan? (\textbf{\ml{പദ്ധതി}}, \emph{paddhati}.)
\end{tcolorbox}
\section[(dialogue)]{Why and the house — a second pass}

\label{lesson:ML-R150-second-pass-why-and-the-house}



\begin{quote}
Say the words for a need and a decision.
\end{quote}

\subsection*{Your turn}
Say these aloud:

\begin{itemize}
\raggedright
  \item a need, a decision, a bed, a mirror, a comb
  \item soap, a towel, a bag, a wall, a roof
  \item a step, a well, a lake, an island, a shore
  \item a garden, a leaf, a mango, a banana, grapes
  \item a potato, a tomato, an aubergine, black pepper, garlic
\end{itemize}

\begin{tcolorbox}[breakable,colback=teal!4,colframe=teal!35!black,title={Before you move on}]
A need? (\textbf{\ml{ആവശ്യം}}, \emph{āvaśyaṁ}.) A decision? (\textbf{\ml{തീരുമാനം}}, \emph{tīrumānaṁ}.) A bed? (\textbf{\ml{കട്ടിൽ}}, \emph{kaṭṭil}.) A mirror? (\textbf{\ml{കണ്ണാടി}}, \emph{kaṇṇāṭi}.) A comb? (\textbf{\ml{ചീപ്പ്}}, \emph{cīppŭ}.) Soap? (\textbf{\ml{സോപ്പ്}}, \emph{sōppŭ}.) A towel? (\textbf{\ml{തോർത്ത്}}, \emph{tōrttŭ}.) A bag? (\textbf{\ml{ബാഗ്}}, \emph{bāgŭ}.) A wall? (\textbf{\ml{ഭിത്തി}}, \emph{bhitti}.) A roof? (\textbf{\ml{മേൽക്കൂര}}, \emph{mēlkkūra}.) A step? (\textbf{\ml{പടി}}, \emph{paṭi}.) A well? (\textbf{\ml{കിണറ്}}, \emph{kiṇaṟŭ}.) A lake? (\textbf{\ml{തടാകം}}, \emph{taṭākaṁ}.) An island? (\textbf{\ml{ദ്വീപ്}}, \emph{dvīpŭ}.) A shore? (\textbf{\ml{തീരം}}, \emph{tīraṁ}.) A garden? (\textbf{\ml{പൂന്തോട്ടം}}, \emph{pūntōṭṭaṁ}.) A leaf? (\textbf{\ml{ഇല}}, \emph{ila}.) A mango? (\textbf{\ml{മാമ്പഴം}}, \emph{māmbaḻaṁ}.) A banana? (\textbf{\ml{വാഴപ്പഴം}}, \emph{vāḻappaḻaṁ}.) Grapes? (\textbf{\ml{മുന്തിരി}}, \emph{muntiri}.) A potato? (\textbf{\ml{ഉരുളക്കിഴങ്ങ്}}, \emph{urulakkiḻaṅṅŭ}.) A tomato? (\textbf{\ml{തക്കാളി}}, \emph{takkāḷi}.) An aubergine? (\textbf{\ml{വഴുതന}}, \emph{vaḻutana}.) Black pepper? (\textbf{\ml{കുരുമുളക്}}, \emph{kurumuḷakŭ}.) Garlic? (\textbf{\ml{വെളുത്തുള്ളി}}, \emph{veḷuttuḷḷi}.)
\end{tcolorbox}
